\chapter{Rarefied Gas Dynamics and Nonequilibrium Closure}\label{ch:17}
\begin{chaptersummary}
Velocity moments connected kinetics to continuum models in Chapters 1 and 2. Stronger nonequilibrium calls for reconsidering the retained state and closure. Knudsen-number comparisons organize Chapman--Enskog, Grad, and simplified collision models, followed by DSMC's stochastic-particle approximation of Boltzmann dynamics. Model, closure, and sampling errors can then be tracked separately.
\end{chaptersummary}
\learninggoals{Distinguish Chapman--Enskog and finite-moment closure; identify the variables in Grad 13-moment and ES--BGK models; assess continuum validity using local scales; identify DSMC sampling error.}
\studyroute{Hyperbolicity proofs for higher-moment systems, general H theorems, and full DSMC collision-sampling constructions are omitted. Transport coefficients depend on the chosen collision kernel.}{Grad 13 moment; Chapman--Enskog; ES BGK; Knudsen layer; DSMC}{For planar Couette flow, compare Newtonian stress, Grad stress, and the kinetic distribution. Identify what boundary data and bulk closure each omit in a Knudsen layer.}

\section{Knudsen Number and Nonequilibrium Information}\index{Rarefied gases}
The ratio $\mathrm{Kn}=\lambda/L$ compares mean free path with a macroscopic length.
Conventional continuum, slip, transition, and free-molecular ranges are modeling guidelines;
local gradients, wall layers, and accuracy targets also determine applicability.
For example, $L_T=T/|\nabla T|$ gives a temperature-gradient Knudsen number $\lambda/L_T$.
Even for small global $\lambda/L$, shocks and wall regions may require nonequilibrium descriptions.

\section{Two Closure Routes: Asymptotic Manifolds and Moment Systems}
At small Knudsen number, the Chapman--Enskog expansion is
\[
f=M[n,\bm u,T]+\mathrm{Kn}\,f^{(1)}
+\mathrm{Kn}^2f^{(2)}+\cdots.
\]
Nonequilibrium corrections have zero conserved moments, so $n,\bm u,T$ are defined by Maxwellian parameters.
Velocity integration of the first correction gives Newtonian stress and Fourier heat flux;
their transport coefficients depend on the linearized collision operator.

Grad methods instead retain higher moments as independent evolving variables.
For a three-dimensional monatomic gas, the 13 moments are
$n$, $\bm u$, $T$, five traceless-stress components, and three heat-flux components.
A corresponding distribution representation has the structure
\[
f_{13}=M\left[1+
A_{ij}\left(c_ic_j-\tfrac13|\bm c|^2\delta_{ij}\right)
+B_i c_i(|\bm c|^2-5\theta)\right],
\qquad \theta=k_BT/m.
\]
Here $A_{ij}$ is symmetric and traceless, with $A,B$ determined by the prescribed stress and heat flux.
Substitution into higher moments gives a finite-moment closure~\cite{grad1949}.
Polynomial corrections can become negative in velocity tails, so realizability and hyperbolicity require explicit checks;
regularized moment models must be selected for their stated regimes.

\conceptbridge{Why do asymptotic and Grad closures differ if both describe stress and heat flux?}
{An asymptotic closure uses small Knudsen number to constrain the nonequilibrium distribution, expressing stress and heat flux through macroscopic gradients.
Grad models retain these moments as independent variables and close the higher moments in their evolution equations.
They therefore use different state spaces and time scales, with different errors and validity regions.}
{Small-parameter asymptotic relations / independent higher moments $\longrightarrow$ two closure families.}

\section{BGK and ES--BGK}
Continuous BGK relaxes toward a Maxwellian while preserving mass, momentum, and energy;
its single relaxation time links transport processes, giving $\mathrm{Pr}=1$ for the standard monatomic model.
ES--BGK uses an anisotropic Gaussian:
\[
G_\chi=\frac{n}{(2\pi)^{3/2}\sqrt{\det\mathcal T_\chi}}
\exp\left(-\frac12\bm c^{\mathsf T}\mathcal T_\chi^{-1}\bm c\right),
\]
\[
\mathcal T_\chi=(1-\chi)\theta I+\chi\Theta,\qquad
\Theta=\bm P/\rho,\qquad
C[f]=\frac{G_\chi-f}{\tau_c}.
\]
The identity $\tr\Theta=3\theta$ preserves energy; the covariance must remain positive definite.
This parameterization allows $\mathrm{Pr}=1/(1-\chi)$;
a common monatomic choice is $\chi=-1/2$, giving $\mathrm{Pr}=2/3$,
with viscosity matched to the model's collision-frequency definition~\cite{holway1966}.
Changes in discrete velocities, quadrature, or time stepping require renewed checks of collision invariants and entropy.

\section{DSMC and Error Separation}
DSMC represents a distribution with weighted simulation particles, alternating free streaming and random collisions in small cells.
Collision partners and scattering angles are sampled from the Boltzmann kernel;
wall or inlet distributions generate incoming particles~\cite{bird1994}.
Finite particle counts introduce statistical fluctuations; cell size and time step introduce numerical bias.
Resolving mean free paths, collision times, and local gradients guides refinement.

To compare continuum and kinetic models, first fix collision physics,
then reduce mesh, time-step, and sampling errors separately.
This separates closure error from numerical error.

\aiexercise{Multiply a standard normal density by $1+\alpha H_3(v)$, with $H_3(v)=v^3-3v$. Explain why any nonzero $\alpha$ can produce negative velocity tails. Calculate zeroth, first, and second moments, and distinguish low-order moment matching from positivity.}
{Design Couette-flow comparisons varying Knudsen number, velocity-space resolution, and DSMC sample count separately. Measure slip velocity and stress, distinguishing closure, boundary-model, and sampling errors.}
